\chapter{System Design}
\label{chap:system-design}

\section{Design Goals}
\label{sec:design-goals}

The framework is designed around four guiding principles: \emph{reproducibility},
\emph{extensibility}, \emph{fairness}, and \emph{traceability}. Together these principles
constrain almost every interface, data structure and naming convention in the codebase.

\paragraph{Reproducibility.}
Every benchmark run is fully described by a single serialisable workload
configuration object. The \code{config_hash()} method serialises this object into
canonically sorted JSON and computes a truncated (8-character) SHA-256 digest. The
hash is stable across machines and Python versions and is used to group repeated
runs of the same workload specification for statistical comparison.

The \texttt{config\_hash()} is deliberately a workload hash rather than a full
execution-context hash. Backend choice, AD mode, requested thread count,
observed thread count, OpenMP/Rayon environment variables, software versions and
hardware metadata are stored in separate database fields. The effective
reproducibility key for a timed result is therefore the combination of
\texttt{config\_hash}, engine, AD mode, parallelism metadata and environment
metadata, not \texttt{config\_hash} alone.

All sources of randomness are explicitly seeded. For a given engine, machine,
thread setting and software environment, the same effective configuration
produces repeatable prices. Different engines use different RNG implementations
(NumPy \texttt{default\_rng}, JAX \texttt{PRNGKey}, C++ \texttt{mt19937\_64},
Rust \texttt{SmallRng}, CUDA \texttt{curand}), so the same seed produces
statistically equivalent but numerically distinct results across engines. In
addition, the C++ and Rust engines use per-thread RNG streams, so changing the
thread count changes the path-to-RNG-stream assignment and can change bit-level
outputs even when the workload hash is unchanged. Cross-engine validation is
therefore performed against analytical references and tolerances rather than
against bit-equal outputs.

\paragraph{Extensibility.}
The framework is organised around two minimal abstractions:
\texttt{WorkloadConfig} (what to simulate) and \texttt{MonteCarloEngine}
(how to simulate it). Adding a workload requires only declaring a new
\texttt{@workload}-decorated dataclass. Adding an engine requires only subclassing
\texttt{MonteCarloEngine} and implementing \texttt{run()}, \texttt{supports()} and
\texttt{supported\_ad\_modes()}. A central \texttt{WORKLOAD\_REGISTRY} maps string
identifiers to workload constructors and is consumed by both the REST API and the
React frontend, so new workloads appear automatically in the UI without
backend--frontend duplication.

\paragraph{Fairness.}
All engines are evaluated under a single benchmarking protocol implemented in
\texttt{BenchmarkRunner}. Each trial consists of a configurable warmup phase whose
results are discarded, followed by a fixed number of timed repetitions. Warmup is
essential for JIT-compiled backends: a cold JAX call at $M=10{,}000$ takes
approximately $115$~ms, compared to approximately $0.6$~ms once XLA has cached the
compiled binary. Excluding the cold compile cost from timing isolates the
steady-state computational performance that is relevant in a production setting.

\paragraph{Traceability.}
Every row in the SQLite database carries the full serialised configuration
(\texttt{config\_json}), the analytical reference values, all timing statistics,
and a complete snapshot of the execution environment (Python version, platform,
CPU model and core count, RAM, NumPy version, JAX version). Any historical run
can be reconstructed exactly, even if the corresponding \texttt{WorkloadConfig}
class is later modified.

\section{High-Level Architecture}
\label{sec:architecture}

The framework follows a strict four-layer architecture, with two additional layers
providing the REST API and the React frontend. Each layer depends only on the
layers below it, so any component can be replaced or extended in isolation.

\begin{figure}[h]
\centering
\fbox{
\parbox{0.9\textwidth}{
\textbf{System Architecture Schematic}

\vspace{0.5em}
The architecture is composed of six vertically stacked layers:

\begin{enumerate}
    \item \textbf{Frontend} (React + TypeScript + Material UI, Vite, port 5173).
    \item \textbf{REST API} (Flask, port 5050) with a background worker thread.
    \item \textbf{Experiment scripts} (\texttt{run\_european\_baseline.py},
          \texttt{run\_european\_ad\_analysis.py},
          \texttt{run\_european\_scalability.py}) that bypass the API and call
          the runner directly.
    \item \textbf{BenchmarkRunner} (warmup, timed repetitions, AD overhead,
          memory and metadata capture).
    \item \textbf{Engine layer}, horizontally split into five variants:
          \begin{itemize}
              \item \texttt{CPUMonteCarloEngine} (NumPy)
              \item \texttt{JAXMonteCarloEngine} (XLA, JIT, AD)
              \item \texttt{CPPMonteCarloEngine} (OpenMP, pybind11)
              \item \texttt{RustMonteCarloEngine} (Rayon, PyO3)
              \item \texttt{CUDAMonteCarloEngine} (PyCUDA, curand)
          \end{itemize}
    \item \textbf{BenchmarkDB} (SQLite with WAL, single \texttt{runs} table).
\end{enumerate}

Configuration objects (\texttt{WorkloadConfig}) flow downwards from layers 1--3
into the runner; engines return \texttt{(price, greeks)} tuples to the runner;
the runner writes fully-populated \texttt{BenchmarkResult} rows into the
database; the API serves both writes (new runs) and reads (history, summary,
compare-matrix) for the frontend.
}
}
\caption{High-level system architecture and data flow.}
\label{fig:arch}
\end{figure}

\section{Layer~1 --- Workload Configuration}
\label{sec:workload-layer}

The workload layer uses a registry pattern implemented via a single Python
decorator. The \texttt{@workload(name)} decorator simultaneously applies
\texttt{@dataclass} to the class, sets a class-level \texttt{workload\_type}
attribute, and registers the class in a global \texttt{WORKLOAD\_REGISTRY}
dictionary. Adding a new workload therefore requires only writing a new class:
no boilerplate, no changes to the runner, storage, API or frontend.

The base class \texttt{WorkloadConfig} provides four concrete inherited methods
used by every workload:

\begin{itemize}
    \item \texttt{validate()} --- a no-op by default, overridden for cross-field
          constraints (e.g.\ enforcing $\sigma > 0$, $T > 0$, valid
          \texttt{option\_type}).
    \item \texttt{to\_dict()} --- full serialisation via
          \texttt{dataclasses.asdict()} with \texttt{workload\_type} injection.
    \item \texttt{from\_dict()} --- safe deserialisation, ignoring unknown keys
          so that schema migrations do not break older saved configurations.
    \item \texttt{config\_hash()} --- the 8-character SHA-256 fingerprint of the
          sorted serialised configuration.
\end{itemize}

There are no abstract methods on \texttt{WorkloadConfig}; every subclass inherits
working implementations. The free function \texttt{config\_from\_dict()} is the
single deserialisation entry point: it reads the \texttt{workload\_type} key,
looks up the corresponding class in \texttt{WORKLOAD\_REGISTRY}, and delegates
to that class's \texttt{from\_dict()}.

\subsection{Workload~1 --- European Option}
\label{sec:workload-european}

\texttt{EuropeanOptionConfig} (\texttt{workload\_type = "european"}) describes a
plain vanilla European call or put under constant-volatility geometric Brownian
motion (GBM). Its fields are summarised in Table~\ref{tab:european-fields}.

\begin{table}[h]
\centering
\caption{Fields of \texttt{EuropeanOptionConfig}.}
\label{tab:european-fields}
\begin{tabular}{lll}
\hline
Field & Default & Description \\
\hline
\texttt{S0} & 100.0 & initial spot price \\
\texttt{K} & 100.0 & strike price \\
\texttt{r} & 0.05 & continuously-compounded risk-free rate \\
\texttt{sigma} & 0.20 & annualised volatility \\
\texttt{T} & 1.0 & time to maturity (years) \\
\texttt{option\_type} & \texttt{"call"} & \texttt{"call"} or \texttt{"put"} \\
\texttt{N} & 1 & number of time steps (single-step exact) \\
\texttt{M} & 10\,000 & number of Monte Carlo paths \\
\texttt{seed} & 42 & RNG seed for reproducibility \\
\hline
\end{tabular}
\end{table}

The pricing model uses single-step exact GBM. The terminal stock price is
\[
S_T \;=\; S_0\,\exp\!\left(\left(r - \tfrac{1}{2}\sigma^2\right)T \;+\; \sigma\sqrt{T}\,Z\right),
\qquad Z \sim \mathcal{N}(0,1),
\]
and the option price is the discounted expected payoff. This is the reference
workload for all engine and AD comparisons because Black--Scholes provides an
exact analytical benchmark for both prices and Greeks.

Analytical validation is built directly into the framework. The
\texttt{CPUMonteCarloEngine} module exposes Black--Scholes closed-form
functions for the call/put price, Delta, Vega and Rho:
\[
\Delta_{\text{call}} = \Phi(d_1), \qquad
\mathrm{Vega} = S_0\,\phi(d_1)\sqrt{T}, \qquad
\rho_{\text{call}} = K\,T\,e^{-rT}\Phi(d_2),
\]
where $d_1 = \bigl(\ln(S_0/K) + (r+\tfrac{1}{2}\sigma^2)T\bigr) / (\sigma\sqrt{T})$
and $d_2 = d_1 - \sigma\sqrt{T}$. These values are computed and stored alongside
every Monte Carlo result, enabling automatic absolute and relative error metrics.

\subsection{Workload~2 --- European Local-Volatility Option}
\label{sec:workload-localvol}

\texttt{EuropeanLocalVolConfig} (\texttt{workload\_type = "european\_local\_vol"})
extends the European workload to a multi-step local-volatility model. The
configuration adds the parameters in Table~\ref{tab:localvol-fields}.

\begin{table}[h]
\centering
\caption{Additional fields of \texttt{EuropeanLocalVolConfig}.}
\label{tab:localvol-fields}
\begin{tabular}{lll}
\hline
Field & Default & Description \\
\hline
\texttt{M} & 100\,000 & paths (larger than European for multi-step precision) \\
\texttt{N} & 252 & Euler steps (daily over one year) \\
\texttt{sigma\_min} & 0.01 & volatility floor \\
\texttt{theta} & auto & polynomial parameters $[a_0, a_1, a_2, b_1]$ \\
\hline
\end{tabular}
\end{table}

The default $\theta$ is set in \texttt{\_\_post\_init\_\_} so that
$\sigma_{\min} + \mathrm{softplus}(a_0) = 0.20$, i.e.\ the model reduces to a
$20\%$ flat-volatility GBM. The closed-form solution is
$a_0 = \log(\mathrm{expm1}(0.20 - \sigma_{\min}))$. When
$a_1 = a_2 = b_1 = 0$, the model collapses to constant-volatility GBM and the
Monte Carlo price converges to the Black--Scholes price as $M$ and $N$ increase.
This provides an automatic regression test embedded in the default configuration.

The dynamics are
\[
dS \;=\; r\,S\,dt \;+\; \sigma(S, t;\theta)\,S\,dW,
\]
discretised in log-space using a log-Euler scheme:
\[
S_{n+1} \;=\; S_n \exp\!\left(\left(r - \tfrac{1}{2}\sigma_n^2\right)\Delta t \;+\; \sigma_n\sqrt{\Delta t}\,Z_n\right).
\]
At each step the local volatility is computed from a parametric monomial surface:
\[
x_n = \ln(S_n / S_0), \qquad
\mathrm{raw}_n = a_0 + a_1 x_n + a_2 x_n^2 + b_1 t_n, \qquad
\sigma_n = \sigma_{\min} + \mathrm{softplus}(\mathrm{raw}_n).
\]
The softplus activation is implemented in a numerically stable form
$\max(\mathrm{raw},0) + \log(1 + \exp(-|\mathrm{raw}|))$, which avoids overflow
for large positive inputs and catastrophic cancellation for large negative
inputs. This formulation is reproduced \emph{identically} in NumPy, JAX, C++
and Rust so that the engines are differentiated only by their execution model,
not by numerical behaviour.

The AD targets for this workload are Delta (\(\partial P / \partial S_0\)) plus
the full $\theta$-gradient \(\partial P / \partial a_0\), \(\partial P / \partial a_1\),
\(\partial P / \partial a_2\), \(\partial P / \partial b_1\),
\(\partial P / \partial \sigma_{\min}\) --- six inputs in total, computed in a
single reverse-mode backward pass.

\section{Layer~2 --- Engine Interface}
\label{sec:engine-layer}

The \texttt{MonteCarloEngine} abstract base class defines a minimal interface:

\begin{verbatim}
class MonteCarloEngine:
    def run(self, config, ad_mode) -> tuple[float, dict | None]
    def supports(self, workload_type: str) -> bool
    def supported_ad_modes(self) -> tuple[str, ...]
\end{verbatim}

The \texttt{ad\_mode} argument takes one of three string values:
\texttt{"none"}, \texttt{"forward"} or \texttt{"reverse"}. When \texttt{ad\_mode}
is \texttt{"none"} or the engine does not implement AD for the given workload,
the returned \texttt{greeks} dictionary is \texttt{None}. The runner inspects
\texttt{supports()} and \texttt{supported\_ad\_modes()} before dispatching a
configuration, so engines that cannot service a request are silently skipped
rather than raising errors.

\paragraph{Engine~1 --- NumPy / CPU.}
\texttt{CPUMonteCarloEngine} is the single-threaded NumPy reference
implementation. It supports the implemented non-AD workloads and
\texttt{ad\_mode = "none"} only.
For the European workload it uses \texttt{np.random.default\_rng(seed)} to
generate all $M$ standard normals as a single array and evaluates the terminal
price and payoff vectorially --- no Python-level loop over paths. For the
local-volatility workload it implements the $N$-step log-Euler scheme in
Python, with each step operating vectorially on an array of $M$ paths. The
correctness of all other engines is verified against this implementation.

\paragraph{Engine~2 --- JAX / XLA.}
\texttt{JAXMonteCarloEngine} is the only engine that exposes
\(\texttt{ad\_mode} \in \{\texttt{none}, \texttt{forward}, \texttt{reverse}\}\).
The European kernel is a
\emph{module-level} \texttt{@jax.jit} function decorated with
\texttt{functools.partial(jax.jit, static\_argnames=("option\_type",))}. Defining
the kernel at module scope ensures JAX traces and compiles it exactly once per
\texttt{(shape, dtype, option\_type)} signature; the random sample $Z$ is passed
as an argument rather than captured in a closure, so the compiled XLA
computation is reused across every timed repetition. The local-volatility
kernel uses \texttt{jax.lax.scan} to implement the $N$-step Euler loop, with
\texttt{static\_argnames=("option\_type", "N")} so XLA emits specialised code
for each call/put and step-count combination.

For the European workload, reverse-mode AD is implemented as a \emph{single}
\texttt{jax.grad(price\_fn, argnums=(0,1,2))} call returning
$(\partial P/\partial S_0,\; \partial P/\partial r,\; \partial P/\partial \sigma)$
in one backward pass. An earlier version invoked \texttt{jax.grad} three times
with \texttt{argnums=0,1,2}, which was correct but tripled both compute time and
peak memory --- a result that materially distorts the AD overhead measurement.
Forward-mode uses three \texttt{jax.jvp} calls probing the standard basis tangent
vectors $(1,0,0)$, $(0,1,0)$ and $(0,0,1)$, exploiting the fact that forward mode
is cheaper than reverse when the number of inputs is small relative to the
number of outputs.

For the local-volatility workload, reverse-mode uses a single
\code{jax.grad(..., argnums=(0,1,2,3,4,5))} call that differentiates through
the full $N$-step \texttt{scan}; forward-mode uses six \texttt{jax.jvp} calls.
JAX differentiates through \texttt{jax.lax.scan} efficiently via its
reverse-mode accumulation, avoiding the explosion in graph size that would
occur if the loop were unrolled.

\paragraph{Engine~3 --- C++ / OpenMP.}
\texttt{CPPMonteCarloEngine} is a thin Python wrapper around a native
implementation in \texttt{cpu\_engine.cpp} / \texttt{cpu\_engine.hpp}, exposed via
\texttt{pybind11} bindings. Each OpenMP thread owns its own \texttt{std::mt19937\_64}
generator seeded as \texttt{seed + thread\_id}; an explicit \texttt{\#pragma omp
parallel} block with a reduction clause accumulates partial sums without locks.
Results are fully deterministic for a given \texttt{(M, seed, thread\_count)}
triple; changing the thread count changes the per-path RNG assignment and produces
statistically equivalent but numerically different results. The local-vol kernel
reproduces exactly the same log-Euler discretisation and stable softplus formula
as the NumPy and JAX engines.

\paragraph{Engine~4 --- Rust / Rayon.}
\texttt{RustMonteCarloEngine} wraps a native implementation in
\texttt{benchmarking/rust/src/lib.rs}, exposed via PyO3 bindings and built with
\texttt{maturin}. The engine uses \texttt{SmallRng} (Xoshiro128++) from the
\texttt{rand} crate, with each Rayon worker thread seeded as
\texttt{seed XOR thread\_index} to mirror the C++ per-thread offset pattern.
Sampling uses \texttt{rand\_distr::StandardNormal}. The per-path simulation is a
pure scalar loop with no heap allocation; the entire $M$-path simulation is
expressed as a Rayon \texttt{into\_par\_iter} chain with \texttt{map\_init} for
per-thread RNG initialisation and a \texttt{sum} reduction. The softplus
is implemented inline as \texttt{z.max(0.0) + (-z.abs()).exp().ln\_1p()}.

\paragraph{Engine~5 --- CUDA / PyCUDA.}
\texttt{CUDAMonteCarloEngine} compiles and launches CUDA kernels at runtime via
PyCUDA. The engine is probed at server startup with a 128-path trial run; if
CUDA initialisation fails it is silently marked unavailable and excluded from
the engine registry. The kernel uses one thread per Monte Carlo path. Each
thread initialises its own \texttt{curandState} via
\texttt{curand\_init(seed, thread\_index, 0, \&state)}, using the thread index
as the sequence number so that statistical independence is guaranteed
regardless of the block/grid layout. Each thread draws $Z$ via
\texttt{curand\_normal\_double}, computes the GBM terminal price, and writes
the discounted payoff to a global output array which the host reduces with
\texttt{np.mean}. Forward-mode Delta is implemented in a separate kernel using
the pathwise estimator: $\partial S_T / \partial S_0 = S_T/S_0$, so the pathwise
Delta is $\mathbf{1}_{\{S_T > K\}}(S_T/S_0)$ for a call and
$-\mathbf{1}_{\{S_T < K\}}(S_T/S_0)$ for a put. The CUDA engine supports the
European workload only.

\section{Layer~3 --- Benchmark Runner}
\label{sec:runner-layer}

\texttt{BenchmarkRunner} encapsulates the full benchmarking protocol. Given an
engine instance and an optional display name, its \texttt{run()} method executes
the following sequence:

\begin{enumerate}
    \item Call \texttt{config.validate()}; abort early on failure.
    \item Execute \texttt{num\_warmup} warmup runs whose results are discarded;
          this excludes XLA compilation from timed measurements.
    \item Execute \texttt{num\_runs} timed runs, recording individual wall-clock
          times via \texttt{time.perf\_counter()}.
    \item If \texttt{ad\_mode != "none"}, repeat steps 2--3 for a matched
          no-AD baseline (same engine, same config, identical warmup count) and
          compute the AD overhead ratio
          $r_{\text{AD}} = \overline{t}_{\text{AD}} / \overline{t}_{\text{baseline}}$.
    \item Snapshot environment metadata: timestamp, Python version and
          implementation, platform, CPU model and architecture, logical CPU
          count, total RAM in GB, NumPy version, JAX version.
    \item Aggregate timings into mean, standard deviation, min and max
          (in milliseconds), and compute
          \texttt{throughput\_paths\_per\_sec = M / mean\_runtime\_s}.
    \item Return a \texttt{BenchmarkResult} dataclass containing all of the
          above plus the price, the Greeks (if any), the analytical reference
          values, and the absolute and relative errors.
\end{enumerate}

Memory peak is measured via \texttt{psutil}: the process resident-set size is
snapshotted before and after the timed run set, and the peak is reported as
$\max(\mathrm{rss}_{\text{during}}) - \mathrm{rss}_{\text{before}}$. If
\texttt{psutil} is unavailable the field is left blank, degrading gracefully.

\section{Layer~4 --- Storage}
\label{sec:storage-layer}

\code{BenchmarkDB} is a thread-safe SQLite wrapper. It uses a per-call connection
with \code{check_same_thread=False} and an explicit \code{threading.Lock} that
serialises writes. The database is opened in Write-Ahead Logging (WAL) mode, which
allows concurrent reads during writes. This matters because the Flask API serves
read requests from the main thread while the background worker writes results.
The database file lives at \code{results/benchmarks.db} and is created on first
use.

\paragraph{Schema.}
The schema consists of a single \code{runs} table whose columns are defined
declaratively by the \code{_FULL_SCHEMA_COLUMNS} list. Columns are grouped as
follows:

\begin{itemize}
    \item \textbf{Identity:} \code{id} (UUID, primary key), \code{experiment_id},
          \code{experiment_type}, \code{config_hash}.
    \item \textbf{Status lifecycle:} \code{status} (\code{pending}/\code{running}/%
          \code{completed}/\code{failed}), \code{error_message},
          \code{created_at}, \code{started_at}, \code{completed_at}
          (all UTC ISO-8601).
    \item \textbf{Workload definition:} \code{workload_type}, \code{M},
          \code{N}, \code{seed}, \code{config_json}.
    \item \textbf{Engine / implementation:} \code{engine}, \code{language},
          \code{backend}.
    \item \textbf{Parallelism metadata:} \code{num_threads},
          \code{vectorization_flag}, \code{batch_size}.
    \item \textbf{Runtime performance:} \code{mean_runtime_ms},
          \code{std_runtime_ms}, \code{min_runtime_ms},
          \code{max_runtime_ms}, \code{throughput_paths_per_sec}.
    \item \textbf{AD metadata:} \code{ad_mode}, \code{baseline_mean_ms},
          \code{ad_overhead_ratio}.
    \item \textbf{Numerical results:} \code{result_value}, \code{greek_delta},
          \code{greek_vega}, \code{greek_rho}.
    \item \textbf{Analytical reference:} \code{analytical_price},
          \code{analytical_delta}, \code{analytical_vega},
          \code{analytical_rho}.
    \item \textbf{Error metrics:} \code{abs_price_error}, \code{rel_price_error}
          and matching columns for Delta, Vega and Rho.
    \item \textbf{Resources:} \code{memory_peak_mb}.
    \item \textbf{Environment:} \code{cpu_model}, \code{cpu_architecture},
          \code{cpu_count}, \code{memory_gb}, \code{platform},
          \code{python_version}, \code{numpy_version},
          \code{jax_version}, \code{blas_backend}.
    \item \textbf{Cloud and profiling:} \code{cloud_provider},
          \code{instance_type}, \code{cloud_region},
          \code{cloud_zone}, \code{hourly_price_usd},
          \code{cost_per_run}, profiler-selection fields, and
          SHA promotion metadata where available.
\end{itemize}

\paragraph{Auto-migration.}
On startup, \texttt{\_init\_schema()} compares the existing table columns
against the authoritative list using \texttt{PRAGMA table\_info} and issues an
\texttt{ALTER TABLE ADD COLUMN} for any missing column. Existing databases
therefore continue to work unmodified as the schema evolves --- a property that
has proved valuable while iterating on the AD and cloud columns during
development.

\paragraph{Write paths.}
Experiment scripts use \texttt{store\_run\_full()}, which inserts a fully-populated
\texttt{completed} row in a single transaction, bypassing the pending/running
state machine. The API uses the
\texttt{create\_run()} $\rightarrow$ \texttt{mark\_running()} $\rightarrow$
\texttt{mark\_completed()} lifecycle to support asynchronous execution.

\section{Layer~5 --- REST API}
\label{sec:api-layer}

The Flask server (\texttt{benchmarking/api/server.py}, port 5050) is intentionally
small. Flask-CORS is enabled for cross-origin requests from the frontend. A single
daemon background thread polls \texttt{BenchmarkDB} for pending runs and executes
them sequentially via \texttt{BenchmarkRunner}, transitioning the row through the
\texttt{running} and \texttt{completed} (or \texttt{failed}) states. The endpoints
are summarised in Table~\ref{tab:api-endpoints}.

\begin{table}[h]
\centering
\caption{REST API endpoints.}
\label{tab:api-endpoints}
\begin{tabular}{lll}
\hline
Endpoint & Method & Purpose \\
\hline
\texttt{/api/runs} & POST & submit a new benchmark run \\
\texttt{/api/runs} & GET & list runs (filter by workload/engine/status) \\
\texttt{/api/runs/:id} & GET & poll a single run by UUID \\
\texttt{/api/engines} & GET & list available engines, workloads and AD modes \\
\texttt{/api/workloads} & GET & list registered workload types \\
\texttt{/api/summary} & GET & aggregate counts, fastest run, breakdowns \\
\texttt{/api/compare\_matrix} & GET & matrix aggregated by (engine, ad\_mode, config) \\
\hline
\end{tabular}
\end{table}

Engine availability is determined at startup. The C++ and CUDA engines may not
be present in every environment; they are excluded from \texttt{/api/engines}
without erroring if their initialisation probe fails.

\section{Layer~6 --- Frontend}
\label{sec:frontend-layer}

The frontend is a React~+~TypeScript single-page application built with Vite
and styled with Material UI, listening on port 5173. It communicates with the
Flask backend through a typed HTTP client in \texttt{src/api/client.ts}. The
application has five pages:

\begin{itemize}
    \item \textbf{SimulatePage} --- submit a new run and poll for completion.
          The \texttt{SimulationForm} component is generated dynamically from
          the workload \texttt{SCHEMA} returned by the API, so new workloads
          appear automatically without frontend changes.
    \item \textbf{ComparePage} --- engine comparison matrix backed by
          \texttt{/api/compare\_matrix}.
    \item \textbf{HistoryPage} --- filterable table of every run.
    \item \textbf{SummaryPage} --- aggregate dashboard.
    \item \textbf{AnalysisPage} --- AD overhead and Greek-accuracy analysis.
\end{itemize}

\section{Cloud Cost-Performance Design}
\label{sec:cloud-cost-design}

Cloud resource selection is treated as a practical systems question rather
than a purely local benchmark. The final design therefore extends each
benchmark row with cloud metadata whenever the run is
executed on a cloud instance. The relevant fields record the provider, instance
type, region, zone, hourly price assumption, git commit, and derived cost per
run. This keeps cost-performance analysis tied to the exact environment in
which a run was measured.

The cost model is intentionally simple:
\[
    \text{cost per run}
    =
    \frac{\text{runtime seconds}}{3600}
    \times
    \text{hourly instance price}.
\]
This does not attempt to model billing granularity, storage, data transfer, or
queueing costs. Its purpose is to compare benchmark configurations under a
common pricing assumption. This is sufficient for identifying whether a faster
instance is also cheaper per run for the measured workload, while avoiding
unsupported claims about long-term cloud bills or all future machine prices.

Cloud fields are stored alongside the same workload, engine, timing, numerical,
and environment metadata as local runs. As a result, cloud and local experiments
share the same analysis path: report artefacts can be regenerated from the
database without manually copying console output into tables.

\section{Budget-Aware Profiler Design}
\label{sec:profiler-design}

Full-grid benchmarking gives the cleanest reference point, but it scales poorly
with the number of workloads, engines, AD modes, path counts, and instance
types. The framework therefore includes profiler logic that uses cheaper probe
runs to select a subset of configurations for full-scale benchmarking. The
important design principle is that the profiler is not trusted by construction:
it is evaluated against a full-grid baseline using explicit selection-quality
metrics.

\subsection{Probe-Based Baseline}
\label{sec:old-profiler-design}

The first profiler uses low-budget probes to rank all candidate configurations
and then selects a conservative subset for full benchmarking. It includes
diversity safeguards so that selection is not determined only by a single global
top-\(k\) ranking. This makes it a useful baseline: it can save full runs, but
it still evaluates a relatively large selected subset to protect against probe
noise and workload-specific crossover behaviour.

\subsection{Successive Halving Profiler}
\label{sec:sha-profiler-design}

The successive-halving (SHA) profiler evaluates all candidate configurations at
the cheapest probe budget, promotes a smaller active set to a larger probe
budget, and repeats this process before choosing the final full-scale runs. In
the final cloud experiment, the budget is the number of Monte Carlo paths, so
the profiler asks whether small-path measurements are sufficiently predictive of
full-path performance.

The design deliberately records intermediate SHA state rather than only the
final selected set. For each promotion stage, the database or generated
artefacts capture the active configuration count and probe budget. These records
allow the report to show how many candidates are discarded at each stage and to
measure the probe path budget consumed by the profiler.

Selection quality is measured using:
\begin{itemize}
    \item \textbf{full runs saved}: the reduction in expensive full-scale
          benchmark evaluations relative to the full grid;
    \item \textbf{Pareto recovery}: the fraction of full-grid non-dominated
          configurations retained by the profiler;
    \item \textbf{runtime regret}: the percentage loss relative to the fastest
          full-grid configuration;
    \item \textbf{cost regret}: the percentage loss relative to the cheapest
          full-grid configuration; and
    \item \textbf{rank correlation}: Spearman agreement between probe rankings
          and full-scale rankings.
\end{itemize}

These metrics make the profiler result falsifiable. A profiler that saves runs
but misses the Pareto frontier or introduces high regret would not be considered
successful. Conversely, a profiler that preserves the full-grid decisions while
reducing full runs contributes directly to the cloud resource-selection goal of
the project.

\section{Design Decisions and Trade-offs}
\label{sec:design-decisions}

\paragraph{Decorator-based workload registration.}
An earlier design made \texttt{WorkloadConfig} an abstract base class with
abstract properties (\texttt{workload\_type}, \texttt{M}, \texttt{seed}) and
abstract methods (\texttt{validate()}, \texttt{to\_dict()}, \texttt{from\_dict()}).
That design forced every workload to re-implement the same boilerplate and
produced a class of bugs in which abstract method implementations drifted out
of sync. The refactored design replaces all of this with a single
\texttt{@workload(name)} decorator that wires up dataclass conversion, sets
\texttt{workload\_type}, and registers the class. New workloads need only declare
their fields and optionally override \texttt{validate()}.

\paragraph{Module-level JAX kernels.}
Placing JAX kernels at module level (rather than as instance methods or nested
functions) is critical for compilation caching. JAX traces and compiles based
on function identity and on argument shapes and dtypes. If the kernel were a
nested function inside \texttt{run()}, each call to \texttt{run()} would create
a fresh function object and defeat the JIT cache, regenerating the XLA binary
on every timed repetition.

\paragraph{Single-pass reverse-mode AD.}
Using a single \texttt{jax.grad(price\_fn, argnums=(0,1,2))} call rather than
three separate \texttt{jax.grad} calls is functionally equivalent but materially
changes the measured AD overhead. The single-pass form shares the backward
graph across all output partial derivatives and is consistent with how AD is
deployed in production quant systems; this is the form benchmarked in this
project.

\paragraph{\texttt{jax.lax.scan} for the local-volatility loop.}
The $N$-step Euler loop must be expressed as \texttt{jax.lax.scan} rather than
as a Python \texttt{for} loop. A Python loop would cause JAX to unroll the
entire computation during tracing, producing an XLA graph with $N=252$ copies
of the step body --- both compile time and memory consumption grow with $N$.
\texttt{jax.lax.scan} compiles to a single loop in the XLA graph and lets the
compiler treat it as a unit.

\paragraph{SQLite over PostgreSQL.}
SQLite with WAL is sufficient for the expected single-machine workload (up to a
few hundred runs per day). It requires no external service, ships as part of
Python, and is straightforward to back up by copying a single file. WAL allows
concurrent reads during writes, which is the only concurrency requirement
imposed by the architecture. If the framework were scaled beyond a
single-machine execution model, PostgreSQL would be a natural replacement.

\paragraph{Per-thread RNG streams.}
The C++ and Rust engines seed each worker thread independently. This sacrifices
bitwise reproducibility across thread counts in exchange for lock-free parallel
sampling. The framework records \texttt{num\_threads} in the database so that
runs at different thread counts are treated as statistically distinct execution
settings rather than as bitwise repeats of the same run.

\section{Summary}
\label{sec:design-summary}

The framework consists of six clearly separated layers: workload configuration,
engine, runner, storage, API and frontend. The four implementation layers are
strictly decoupled through two minimal abstractions
(\texttt{WorkloadConfig} and \texttt{MonteCarloEngine}) and a registry-driven
workload model. Every benchmark run is fully described by a serialisable
configuration, fingerprinted by a deterministic hash, and persisted alongside its
analytical reference values and execution environment, ensuring complete
reproducibility and traceability. The next chapter describes how each layer is
implemented in practice.
